\section{Complete finite heat divisors}
\label{sec:finite}

Let $P(t,z)=\sum_{j=0}^N a_j(t)z^{N-j}$ be monic, with $a_0=1$,
and suppose $P_t=-P_{zz}$. This section concerns the complete finite
polynomial, with no omitted exterior. We allow complex coefficients
and roots unless a reality condition is explicitly stated.

\begin{theorem}[Coefficient integrals and a polynomial clock]
\label{thm:finite}
The coefficients satisfy
\begin{equation}
 a_j'=-(N-j+2)(N-j+1)a_{j-2},\qquad a_{-1}=a_{-2}=0.
 \label{eq:coeff-law}
\end{equation}
Consequently
\begin{equation}
 I_j(t)=\sum_{r=0}^{\lfloor j/2\rfloor}
 \frac{t^r}{r!}\frac{(N-j+2r)!}{(N-j)!}a_{j-2r}(t)=a_j(0).
 \label{eq:explicit-integrals}
\end{equation}
For $N\ge2$, set $\theta=-a_2/[N(N-1)]$. Then $\theta'=1$ and
every coefficient of
\begin{equation}
 Z(z)=e^{\theta\partial_z^2}P(t,z)
 \label{eq:autonomous}
\end{equation}
is constant. Its coefficient of $z^{N-2}$ is zero. Its other $N-1$
nonleading coefficients are independent autonomous integrals of the
$N$-dimensional monic coefficient flow.
\end{theorem}
\begin{proof}
Comparing powers in $P_t=-P_{zz}$ gives \eqref{eq:coeff-law}.
The exponential of $\partial_z^2$ terminates on a polynomial and
commutes with differentiation in time. Thus
$\partial_t(e^{t\partial_z^2}P)=0$. Its $j$th coefficient is
\eqref{eq:explicit-integrals}. Similarly,
\[
 \partial_t(e^{\theta\partial_z^2}P)
 =(\theta'-1)\partial_z^2e^{\theta\partial_z^2}P=0.
\]
The $z^{N-2}$ coefficient of $Z$ is $a_2+N(N-1)\theta=0$.
Conversely, an arbitrary monic $Z$ with that coefficient zero,
together with an arbitrary $\theta$, recovers
$P=e^{-\theta\partial_z^2}Z$. Its second coefficient is
$-N(N-1)\theta$, so this is an inverse polynomial coordinate change.
Independence and completeness of the $N-1$ constants follow.
\end{proof}

For simple roots, evaluation of the heat equation at $P(t,x_i(t))=0$
gives
\begin{equation}
 x_i'=\frac{P_{zz}(t,x_i)}{P_z(t,x_i)}
       =2\sum_{j\ne i}\frac1{x_i-x_j}.
 \label{eq:finite-roots}
\end{equation}
Set
\begin{equation}
 c=\frac1N\sum_i x_i,\qquad y_i=x_i-c,\qquad p_k=\sum_i y_i^k.
 \label{eq:center}
\end{equation}
Then $p_0=N$, $p_1=0$, and pairwise cancellation gives $c'=0$.
For any integer $k\ge2$,
\begin{equation}
 p_k'=2k\sum_{i<j}\frac{y_i^{k-1}-y_j^{k-1}}{y_i-y_j}.
 \label{eq:powers}
\end{equation}
In particular,
\begin{equation}
 p_2'=2N(N-1),\qquad p_3'=0,\qquad p_4'=4(2N-3)p_2.
 \label{eq:low-powers}
\end{equation}
For the last identity, the unordered sum of $y_i^2+y_iy_j+y_j^2$
is $(N-1)p_2-p_2/2=(N-3/2)p_2$. This proves every coefficient
in \eqref{eq:low-powers}, including its sign.

\begin{corollary}[Centered homogeneous integrals]
\label{cor:homogeneous}
Let $Q(t,y)=\prod_i(y-y_i)=\sum_{j=0}^N b_jy^{N-j}$, so $b_1=0$
and $b_2=-p_2/2$. Define
\begin{equation}
 \vartheta=\frac{p_2}{2N(N-1)},\qquad
 J_j=[y^{N-j}]\,e^{\vartheta\partial_y^2}Q,
 \quad 3\le j\le N.
 \label{eq:family}
\end{equation}
Each $J_j$ is a translation-invariant symmetric polynomial in the
$x_i$, homogeneous of degree $j$, and a first integral of
\eqref{eq:finite-roots}. The $N-2$ polynomials $J_3,\ldots,J_N$
are independent in the centered coefficient variables. Along with
the centroid they give $N-1$ independent autonomous integrals.
\end{corollary}
\begin{proof}
The center is constant, so $Q_t=-Q_{yy}$ and $\vartheta'=1$.
Theorem~\ref{thm:finite} applies to $Q$, and its inverse coordinate
change now uses $\vartheta,J_3,\ldots,J_N$. Translation leaves every
$y_i$ unchanged. Under $x\mapsto\lambda x$, the centered coefficient
$b_{j-2r}$ has degree $j-2r$ and $\vartheta^r$ has degree $2r$.
Every term in $J_j$ therefore has degree $j$. The coefficient formula
also proves polynomiality and symmetry without selecting root labels.
\end{proof}

Newton identities give $J_3=b_3=-p_3/3$. For $N\ge4$ they give
the useful fourth-order integral
\begin{align}
 C_4:=J_4
 &=\frac{2N-3}{4N(N-1)}p_2^2-\frac14p_4 \label{eq:C4}\\
 &=b_4-\frac{(N-2)(N-3)}{2N(N-1)}b_2^2.
 \notag
\end{align}
Indeed $b_4=p_2^2/8-p_4/4$ and $b_2=-p_2/2$; substitution into
\eqref{eq:family} proves \eqref{eq:C4}. It can also be verified
directly from \eqref{eq:low-powers}. All these statements extend
through collisions because they are polynomial coefficient identities.

\subsection{A collision with positive conserved value}
For $b>0$ and arbitrary real $T$, consider the exact heat polynomial
\begin{equation}
 P(t,z)=(z^2-b^2)^2-(t-T)(12z^2-4b^2)+12(t-T)^2.
 \label{eq:quartic-control}
\end{equation}
Writing $\tau=t-T$, its coefficients are
$a=-2b^2-12\tau$ and $d=b^4+4b^2\tau+12\tau^2$ in
$P=z^4+az^2+d$. Direct differentiation verifies $P_t=-P_{zz}$,
and \eqref{eq:C4} becomes
\begin{equation}
 C_4=d-a^2/12=2b^4/3>0.
 \label{eq:positive-collision}
\end{equation}
Nevertheless $\pm b$ are double zeros at $t=T$.

More generally, for a real centered even quartic $z^4+az^2+d$,
$p_2=-2a$ and $d=C_4+p_2^2/48$. The discriminant in $z^2$ is
\begin{equation}
 \Delta_{z^2}=a^2-4d=p_2^2/6-4C_4.
 \label{eq:quartic-discriminant}
\end{equation}
If $p_2>0$ and $C_4>0$, both $z^2$ roots have positive sum and
positive product. Thus the four $z$ roots are real and simple exactly
when $p_2^2>24C_4$. Equality gives the two nonzero double roots.
This finite threshold uses the variance clock and the invariant jointly.
It contains exactly the information of the direct quartic discriminant;
no same-information advantage is asserted. Positivity of $C_4$ alone
does not exclude a collision.
